
\section{Estimação por máxima verossimilhança}

Considere uma amostra independente
$y_1,\ldots,y_n$ proveniente de uma distribuição
de Poisson inflada em zero (ZIP), com parâmetros
$0<p<1$ e $\lambda>0$.

A função de probabilidade é

\begin{equation}
P(Y=y)=
\begin{cases}
p+(1-p)e^{-\lambda}, & y=0,\\
(1-p)e^{-\lambda}\dfrac{\lambda^y}{y!},
& y>0.
\end{cases}
\end{equation}

Definimos $n_0$ como o número de observações iguais
a zero, $n_+=n-n_0$ como o número de observações
positivas e $S=\sum_{i=1}^{n}y_i$.

A log-verossimilhança, omitindo termos constantes
em relação aos parâmetros, é

\begin{equation}
\ell(p,\lambda)
=
n_0\log\left[p+(1-p)e^{-\lambda}\right]
+n_+\log(1-p)
-n_+\lambda
+S\log\lambda.
\end{equation}

\subsection{Questão 1: Equações de score}

Para obter as condições de primeira ordem da
máxima verossimilhança, calculamos as derivadas
parciais de $\ell(p,\lambda)$.

Definindo

\begin{equation}
A(p,\lambda)=p+(1-p)e^{-\lambda},
\end{equation}

temos

\begin{equation}
\frac{\partial A}{\partial p}
=1-e^{-\lambda},
\qquad
\frac{\partial A}{\partial\lambda}
=-(1-p)e^{-\lambda}.
\end{equation}

Derivando a log-verossimilhança em relação a $p$,

\begin{align}
\frac{\partial\ell}{\partial p}
&=
n_0\frac{1}{A(p,\lambda)}
\frac{\partial A}{\partial p}
-\frac{n_+}{1-p}\\
&=
\frac{n_0(1-e^{-\lambda})}
{p+(1-p)e^{-\lambda}}
-\frac{n_+}{1-p}.
\end{align}

Analogamente, derivando em relação a $\lambda$,

\begin{align}
\frac{\partial\ell}{\partial\lambda}
&=
n_0\frac{1}{A(p,\lambda)}
\frac{\partial A}{\partial\lambda}
-n_++\frac{S}{\lambda}\\
&=
-\frac{n_0(1-p)e^{-\lambda}}
{p+(1-p)e^{-\lambda}}
-n_++\frac{S}{\lambda}.
\end{align}

Assim, as equações de score para um ponto
estacionário interior são

\begin{equation}
\boxed{
\frac{n_0(1-e^{-\lambda})}
{p+(1-p)e^{-\lambda}}
-\frac{n_+}{1-p}=0
}
\end{equation}

e

\begin{equation}
\boxed{
-\frac{n_0(1-p)e^{-\lambda}}
{p+(1-p)e^{-\lambda}}
-n_++\frac{S}{\lambda}=0.
}
\end{equation}

Essas equações caracterizam os pontos estacionários
interiores da log-verossimilhança e constituem a
base para a implementação do método de Newton-Raphson.


\subsection{Estimação numérica com \texttt{optim()}}

A estimação por máxima verossimilhança foi realizada
separadamente para os três perfis de risco,
utilizando a função \texttt{optim()} do R.

Para respeitar o espaço paramétrico do modelo,
adotou-se a reparametrização
\[
p=\frac{1}{1+e^{-\alpha}},
\qquad
\lambda=e^\beta.
\]

Foi empregado o método BFGS, considerando nove
combinações de valores iniciais por perfil.
Para cada perfil, selecionou-se a solução de
maior log-verossimilhança.

\begin{table}[htbp]
\centering
\caption{Estimativas de máxima verossimilhança
do modelo ZIP por perfil de risco.}

\begin{tabular}{lrrr}
\toprule
Perfil & $\hat{p}$ & $\hat{\lambda}$ & Log-verossimilhança\\
\midrule
baixo & 0.5644 & 1.6874 & -316.9396\\
medio & 0.2822 & 3.0127 & -104.7402\\
alto & 0.0811 & 4.9679 & 997.6518\\
\bottomrule
\end{tabular}

\end{table}

Todas as nove inicializações de cada perfil
retornaram código de convergência igual a zero.

Nos perfis baixo e médio, os resultados apresentaram
log-verossimilhanças praticamente idênticas.

No perfil alto, cinco inicializações conduziram
à solução superior, enquanto quatro terminaram
próximas de $p=0$, com log-verossimilhança inferior
em aproximadamente 51 unidades.

Portanto, a utilização de múltiplos pontos
iniciais foi importante para identificar
soluções numericamente distintas e selecionar
o melhor ajuste encontrado.

A log-verossimilhança apresentada omite
o termo constante
$-\sum_{i=1}^{n}\log(y_i!)$,
que não interfere nas estimativas de
máxima verossimilhança.
